% architecture-infographic.tex
% echarts-examples 项目架构信息图
%
% 叙事主线：12 个透明背景 ECharts 场景 → 目录驱动的组织 → Vite 构建 →
%           Playwright 离线渲染与 RGBA 质量门 → PNG 输出
%
% 视觉语法：纵向三段式
%   上段 — 项目身份与核心命题（主角：12 场景网格）
%   中段 — 构建与渲染管线（带状路径 + 往返）
%   下段 — 质量验证与技术栈（小倍数 + 表格）
%
% Render: python3 ~/.agents/skills/create-infographic/scripts/render_tikz.py docs/architecture-infographic.tex --workdir . --out-dir docs --engine xelatex --dpi 300 --svg --strict-warnings --strict-geometry --strict-composition --strict-svg --min-labelled-connectors 0.6 --expect-text "echarts-examples"

\documentclass[border=14pt,varwidth=\maxdimen]{standalone}
\usepackage[UTF8,fontset=none]{ctex}
\IfFontExistsTF{Source Han Serif SC}{
  % LetterSpace=5.0：pdftocairo 导出 SVG 时相邻字形 bbox 亚像素相切，
  % 会被 svg/text-text-overlap 误报；微字距让导出 bbox 脱开（见 svg-linter）
  \setmainfont[BoldFont={Source Han Serif SC Bold},BoldFeatures={LetterSpace=5.0},LetterSpace=5.0]{Source Han Serif SC}
  \setsansfont[BoldFont={Source Han Serif SC Bold},BoldFeatures={LetterSpace=5.0},LetterSpace=5.0]{Source Han Serif SC}
  \setCJKmainfont[BoldFont={Source Han Serif SC Bold},BoldFeatures={LetterSpace=5.0},LetterSpace=5.0]{Source Han Serif SC}
  \setCJKsansfont[BoldFont={Source Han Serif SC Bold},BoldFeatures={LetterSpace=5.0},LetterSpace=5.0]{Source Han Serif SC}
}{
  \PackageWarningNoLine{tikz-infographic}{Source Han Serif SC not found; using Fandol}
  \setCJKmainfont[BoldFont={FandolSong-Bold}]{FandolSong-Regular}
  \setCJKsansfont[BoldFont={FandolHei-Bold}]{FandolHei-Regular}
}
\IfFontExistsTF{Sarasa Mono SC}{
  \setmonofont{Sarasa Mono SC}
  \setCJKmonofont{Sarasa Mono SC}
}{
  \PackageWarningNoLine{tikz-infographic}{Sarasa Mono SC not found; using TeX defaults}
  \setmonofont{Latin Modern Mono}
  \setCJKmonofont{FandolFang-Regular}
}
\usepackage{xcolor}
\usepackage{tikz}
\usetikzlibrary{arrows.meta,backgrounds,calc,fit,positioning,shapes.geometric}
\pgfdeclarelayer{bg layer}
\pgfdeclarelayer{edge layer}
\pgfsetlayers{bg layer,edge layer,main}

% ── Semantic palette ────────────────────────────────────────────────────────
% 深色模式呼应项目本身的透明画布 + 冷色调配色
\definecolor{Paper}{HTML}{F7F4EE}
\definecolor{Ink}{HTML}{0F1B2A}
\definecolor{Slate}{HTML}{4F6378}
\definecolor{Fog}{HTML}{CFDBE5}
\definecolor{Deep}{HTML}{142744}      % 场景卡深色底（呼应项目 backgroundColor 风格）
\definecolor{Ocean}{HTML}{2E6B8A}
\definecolor{Teal}{HTML}{2FA496}       % #54d6c6 降饱和
\definecolor{Mint}{HTML}{D6EEE9}
\definecolor{Coral}{HTML}{E2667B}      % #ff6f91 降饱和
\definecolor{Gold}{HTML}{E8B94F}       % #ffd166 降饱和
\definecolor{Violet}{HTML}{8A7CC9}
\definecolor{Sky}{HTML}{5E8FD6}
\definecolor{Panel}{HTML}{EEF2F6}

\tikzset{
  every node/.style={font=\rmfamily},
  % 主路径色带
  spine/.style={draw=Fog, line width=4mm, line cap=butt, line join=round},
  % 阶段路标（编号 + 名称合一的 pill，避免文字压边）
  % inner xsep 需 > 色带半宽 + 1px 间隙，否则 SVG 审计判定色带压字
  stage pill/.style={
    rounded corners=1.8mm, draw=Ocean, fill=Ocean, text=white,
    inner xsep=4.4mm, inner ysep=1.1mm,
    font=\rmfamily\scriptsize\bfseries
  },
  % 系统 / 模块边界框（仅在确有边界时使用）
  module/.style={
    draw=Ocean!40, fill=white, rounded corners=1.6mm, line width=.7pt,
    minimum width=24mm, minimum height=9mm, align=center,
    font=\rmfamily\footnotesize\bfseries, text=Ink
  },
  % 场景缩略卡（12 个小倍数，每个都代表一个真实 ECharts 场景）
  scene card/.style={
    draw=Deep!25, fill=Deep, rounded corners=1.2mm, line width=.6pt,
    minimum width=30mm, minimum height=15mm, inner sep=0pt
  },
  % 场景卡下方的标题 + 标签（文字全部移出卡片，避免压住缩略图）
  scene caption/.style={
    anchor=north west, align=left, text width=29mm, inner sep=0pt,
    font=\rmfamily\scriptsize, text=Ink
  },
  % 边语法
  call/.style={
    -{Latex[length=2.0mm,width=1.4mm]}, draw=Ocean,
    line width=.9pt, rounded corners=1.6mm
  },
  reply/.style={
    -{Latex[length=2.0mm,width=1.4mm]}, draw=Teal,
    densely dashed, line width=.8pt, rounded corners=1.6mm
  },
  dataflow/.style={
    -{Latex[length=2.0mm,width=1.4mm]}, draw=Gold,
    line width=.9pt, rounded corners=1.6mm
  },
  verify/.style={
    -{Latex[length=2.0mm,width=1.4mm]}, draw=Coral,
    densely dotted, line width=.9pt, rounded corners=2mm
  },
  boundary/.style={draw=Slate!40, dash pattern=on 2pt off 2.2pt, line width=.55pt},
  % 文本风格
  edge label/.style={font=\rmfamily\scriptsize, text=Slate},
  terminal/.style={font=\rmfamily\scriptsize\bfseries, text=Ocean},
  note/.style={font=\rmfamily\scriptsize, text=Slate},
  kicker/.style={font=\rmfamily\scriptsize\bfseries, text=Teal},
  pill/.style={
    rounded corners=2.6mm, fill=Mint, text=Ocean,
    inner xsep=2.6mm, inner ysep=1.0mm,
    font=\rmfamily\scriptsize\bfseries
  },
  section title/.style={font=\rmfamily\large\bfseries, text=Ink},
  section sub/.style={font=\rmfamily\scriptsize, text=Slate},
  family label/.style={font=\rmfamily\footnotesize\bfseries, text=Ocean},
  metric num/.style={font=\rmfamily\Large\bfseries, text=Deep},
  metric label/.style={font=\rmfamily\scriptsize, text=Slate}
}

\begin{document}
\setlength{\parindent}{0pt}
\begin{tikzpicture}[x=1cm,y=1cm]

  % ── 画布底色 ────────────────────────────────────────────────────────────
  \begin{pgfonlayer}{bg layer}
    \fill[Paper,rounded corners=2mm] (-0.6,-2.4) rectangle (19.7,27.9);
  \end{pgfonlayer}

  % ════════════════════════════════════════════════════════════════════════
  %  第一段：项目身份与核心主张
  % ════════════════════════════════════════════════════════════════════════

  \node[kicker,anchor=west] at (0,27.45) {ECharts Examples · 架构信息图};
  \node[anchor=west,font=\rmfamily\Large\bfseries,text=Ink]
    at (0,26.72) {十二个高信息密度场景，一条离线验证管线};
  \node[note,anchor=west] at (0,26.05)
    {透明背景 · Canvas 渲染 · Playwright 截图 · RGBA 质量门};
  \node[pill,anchor=east] at (19.1,26.72) {v0.2.0 · TypeScript + Vite};

  % ── 四个关键指标 ────────────────────────────────────────────────────────
  \node[metric num,anchor=west] at (0,25.08) {12};
  \node[metric label,anchor=north west,align=left] at (0,24.42) {ECharts 场景\\按分析家族分组};

  \node[metric num,anchor=west] at (3.6,25.08) {6};
  \node[metric label,anchor=north west,align=left] at (3.6,24.42) {视觉家族\\流程/层次/网络/时间…};

  \node[metric num,anchor=west] at (7.6,25.08) {1400×900};
  \node[metric label,anchor=north west,align=left] at (7.6,24.42) {每场景画布\\透明 PNG 输出};

  \node[metric num,anchor=west] at (11.6,25.08) {3 重};
  \node[metric label,anchor=north west,align=left] at (11.6,24.42) {质量门\\类型 / 交互 / RGBA};

  \node[metric num,anchor=west] at (15.0,25.08) {100\%};
  \node[metric label,anchor=north west,align=left] at (15.0,24.42) {离线可测\\无网络依赖};

  % 分隔线
  \draw[boundary] (0,23.45) -- (19.1,23.45);

  % ════════════════════════════════════════════════════════════════════════
  %  第二段：12 个场景 · 按分析家族分组
  %  主角区域：视觉上占比最大，直接展示项目的核心产出
  % ════════════════════════════════════════════════════════════════════════

  \node[kicker,anchor=west] at (0,23.05) {场景目录 · Scene Catalog};
  \node[section title,anchor=west] at (0,22.42) {12 个场景，覆盖六大分析家族};
  \node[section sub,anchor=west] at (0,21.9)
    {每个场景都有明确的分析问题、图表家族与复杂度等级，由 catalog.json 统一描述。};

  % catalog.json 引用卡（统一描述全部场景）
  \node[module,anchor=north west,minimum width=50mm,fill=Panel,
        text width=46mm,align=left,inner sep=2mm]
    (catalog) at (13.6,22.9)
    {{\footnotesize\bfseries catalog.json}\\[1pt]
     {\scriptsize 每个场景记录：用途 · 问题 · 家族 · 复杂度 · 标签}};

  % catalog → 场景网格的代表性连接（语义：目录驱动全部 12 个场景）
  \begin{pgfonlayer}{edge layer}
    \draw[call] (catalog.south -| 16.1,0) -- (16.1,20.95);
  \end{pgfonlayer}
  \node[edge label,anchor=west] at (16.35,21.22) {描述全部 12 场景};

  % ── 家族 1：流程 / 漏斗 ────────────────────────────────────────────────
  \node[family label,anchor=west] at (0,21.28) {01 · 流程与转化};

  % Sankey
  \node[scene card,anchor=north west] (sc-sankey) at (0,20.95) {};
  \node[scene caption] at (0.05,19.3)
    {{\bfseries Sankey}\\[0.5pt]\textcolor{Slate}{能源流向}\\\textcolor{Slate}{可再生→储能→需求}};
  % 小装饰：三条互不交叉的流线（填充细带，不用描边——描边会被计为无标签连线）
  \begin{scope}[shift={(sc-sankey.center)},x=1mm,y=1mm]
    \fill[Teal,opacity=.85] (-12,3.64) .. controls (-4,3.64) and (4,2.14) .. (12,2.14)
      -- (12,1.86) .. controls (4,1.86) and (-4,3.36) .. (-12,3.36) -- cycle;
    \fill[Gold,opacity=.85] (-12,0.64) .. controls (-4,0.64) and (4,-0.86) .. (12,-0.86)
      -- (12,-1.14) .. controls (4,-1.14) and (-4,0.36) .. (-12,0.36) -- cycle;
    \fill[Coral,opacity=.85] (-12,-2.36) .. controls (-4,-2.36) and (4,-3.86) .. (12,-3.86)
      -- (12,-4.14) .. controls (4,-4.14) and (-4,-2.64) .. (-12,-2.64) -- cycle;
  \end{scope}

  % Funnel
  \node[scene card,anchor=north west] (sc-funnel) at (3.3,20.95) {};
  \node[scene caption] at (3.35,19.3)
    {{\bfseries Funnel}\\[0.5pt]\textcolor{Slate}{科研招募}\\\textcolor{Slate}{候选人流失节点}};
  % 小装饰：漏斗梯形（填充块之间留缝，无描边）
  \begin{scope}[shift={(sc-funnel.center)},x=1mm,y=1mm]
    \fill[Gold,opacity=.85] (-11,4.5) -- (11,4.5) -- (7,1.0) -- (-7,1.0) -- cycle;
    \fill[Coral,opacity=.85] (-7,0.5) -- (7,0.5) -- (4,-3.0) -- (-4,-3.0) -- cycle;
    \fill[Teal,opacity=.85] (-4,-3.5) -- (4,-3.5) -- (1.8,-6.0) -- (-1.8,-6.0) -- cycle;
  \end{scope}

  % ── 家族 2：层次 / 构成 ────────────────────────────────────────────────
  \node[family label,anchor=west] at (6.8,21.28) {02 · 层次与构成};

  % Sunburst
  \node[scene card,anchor=north west] (sc-sunburst) at (6.8,20.95) {};
  \node[scene caption] at (6.85,19.3)
    {{\bfseries Sunburst}\\[0.5pt]\textcolor{Slate}{研究组合}\\\textcolor{Slate}{三层项目分布}};
  % 小装饰：同心扇环（纯填充、扇区间留角缝，无描边）
  \begin{scope}[shift={(sc-sunburst.center)},x=1mm,y=1mm]
    \fill[Teal,opacity=.85] (0,0) -- (5:2.2) arc (5:105:2.2) -- cycle;
    \fill[Sky,opacity=.85] (0,0) -- (115:2.2) arc (115:225:2.2) -- cycle;
    \fill[Coral,opacity=.85] (0,0) -- (235:2.2) arc (235:345:2.2) -- cycle;
    \foreach \a in {0,45,90,135,180,225,270,315} {
      \fill[Gold,opacity=.55]
        (\a+4:3.4) -- (\a+4:6.6) arc (\a+4:\a+41:6.6) -- (\a+41:3.4) arc (\a+41:\a+4:3.4) -- cycle;
    }
  \end{scope}

  % Treemap
  \node[scene card,anchor=north west] (sc-treemap) at (10.1,20.95) {};
  \node[scene caption] at (10.15,19.3)
    {{\bfseries Treemap}\\[0.5pt]\textcolor{Slate}{资产组合}\\\textcolor{Slate}{价值与风险分布}};
  % 小装饰：嵌套矩形（纯填充、块间留缝）
  \begin{scope}[shift={(sc-treemap.center)},x=1mm,y=1mm]
    \fill[Teal,opacity=.8]   (-11,4.5) rectangle (-3,0.5);
    \fill[Sky,opacity=.8]    (-2.5,4.5) rectangle (11,2.5);
    \fill[Coral,opacity=.8]  (-2.5,2.0) rectangle (5,0.5);
    \fill[Gold,opacity=.8]   (5.5,2.0) rectangle (11,0.5);
    \fill[Gold,opacity=.7]   (-11,0) rectangle (-4,-5.5);
    \fill[Violet,opacity=.8] (-3.5,0) rectangle (2,-5.5);
    \fill[Teal,opacity=.6]   (2.5,0) rectangle (11,-5.5);
  \end{scope}

  % ── 家族 3：网络 / 关系 ────────────────────────────────────────────────
  \node[family label,anchor=west] at (13.6,21.28) {03 · 网络与图};

  % Knowledge Graph
  \node[scene card,anchor=north west] (sc-graph) at (13.6,20.95) {};
  \node[scene caption] at (13.65,19.3)
    {{\bfseries Knowledge Graph}\\[0.5pt]\textcolor{Slate}{气候知识图谱}\\\textcolor{Slate}{社区与桥接节点}};
  % 小装饰：六边形环 + 中心节点（环拆成 6 段填充细矩形、辐条为填充细矩形，
  % 端头缩进藏在节点圆点下；均为填充而非描边，避免被计为无标签连线；
  % 线段不包含圆点中心，避免容器嵌套过深）
  \begin{scope}[shift={(sc-graph.center)},x=1mm,y=1mm]
    \foreach \a in {0,60,120,180,240,300}
      \fill[white!30,rotate around={\a:(0,0)}] (4.95,-1.69) rectangle (5.09,1.69);
    \fill[white!30] (-0.07,1.55) rectangle (0.07,4.65);
    \fill[white!30,rotate around={210:(0,0)}] (-0.07,1.55) rectangle (0.07,4.65);
    \fill[white!30,rotate around={330:(0,0)}] (-0.07,1.55) rectangle (0.07,4.65);
    \foreach \p in {(0,5.8),(-5.0,2.9),(-5.0,-2.9),(0,-5.8),(5.0,-2.9),(5.0,2.9)}
      \fill[Teal] \p circle (1.2);
    \fill[Gold] (0,0) circle (1.5);
  \end{scope}

  % ── 家族 4：时间 / 时序 ────────────────────────────────────────────────
  \node[family label,anchor=west] at (0,18.02) {04 · 时间与模式};

  % Calendar
  \node[scene card,anchor=north west] (sc-calendar) at (0,17.69) {};
  \node[scene caption] at (0.05,16.04)
    {{\bfseries Calendar}\\[0.5pt]\textcolor{Slate}{运营异常日历}\\\textcolor{Slate}{全年日度强度}};
  % 小装饰：日历热力格（纯填充）
  \begin{scope}[shift={(sc-calendar.center)},x=1mm,y=1mm]
    \foreach \i in {0,...,27} {
      \pgfmathtruncatemacro{\cx}{mod(\i,7)}
      \pgfmathtruncatemacro{\cy}{\i/7}
      \pgfmathtruncatemacro{\hash}{mod(\cx*5+\cy*3+7,40)}
      \ifnum\hash<10 \definecolor{tmpc}{HTML}{1a2f4a} \else
      \ifnum\hash<20 \definecolor{tmpc}{HTML}{2A9D8F} \else
      \ifnum\hash<30 \definecolor{tmpc}{HTML}{54d6c6} \else
      \definecolor{tmpc}{HTML}{E8B94F} \fi\fi\fi
      \fill[tmpc] (\cx*3.0 - 9.0, \cy*3.0 - 4.5) rectangle (\cx*3.0 - 6.6, \cy*3.0 - 2.1);
    }
  \end{scope}

  % Theme River
  \node[scene card,anchor=north west] (sc-river) at (3.3,17.69) {};
  \node[scene caption] at (3.35,16.04)
    {{\bfseries Theme River}\\[0.5pt]\textcolor{Slate}{文化主题流}\\\textcolor{Slate}{1960–2031 演变}};
  % 小装饰：主题河面积（纯填充）
  \begin{scope}[shift={(sc-river.center)},x=1mm,y=1mm]
    \fill[Teal,opacity=.7]
      (-12,0) sin (-6,3) cos (0,2) sin (6,4) cos (12,3) --
      (12,0) cos (6,-2) sin (0,-1) cos (-6,-3) sin (-12,0) -- cycle;
    \fill[Gold,opacity=.7]
      (-12,0) sin (-6,-2) cos (0,-1) sin (6,-2) cos (12,0) --
      (12,0) cos (6,-4) sin (0,-3) cos (-6,-5) sin (-12,0) -- cycle;
  \end{scope}

  % Candlestick
  \node[scene card,anchor=north west] (sc-candle) at (6.8,17.69) {};
  \node[scene caption] at (6.85,16.04)
    {{\bfseries Candlestick}\\[0.5pt]\textcolor{Slate}{金融微观结构}\\\textcolor{Slate}{价格/量/波动}};
  % 小装饰：K 线（影线与实体均为填充矩形，无描边穿越）
  \begin{scope}[shift={(sc-candle.center)},x=1mm,y=1mm]
    \fill[Teal]  (-9.25,-1) rectangle (-8.75,4.5); \fill[Teal]  (-9.8,0) rectangle (-8.2,3);
    \fill[Coral] (-6.25,-3) rectangle (-5.75,3.5); \fill[Coral] (-6.8,-1) rectangle (-5.2,2);
    \fill[Teal]  (-3.25,-3.5) rectangle (-2.75,4); \fill[Teal]  (-3.8,-1) rectangle (-2.2,2);
    \fill[Coral] (-0.25,-1) rectangle (0.25,5);    \fill[Coral] (-0.8,1) rectangle (0.8,3);
    \fill[Teal]  (2.75,-1.5) rectangle (3.25,4);   \fill[Teal]  (2.2,1) rectangle (3.8,3);
    \fill[Coral] (5.75,-4) rectangle (6.25,3);     \fill[Coral] (5.2,-2) rectangle (6.8,2);
    \fill[Teal]  (8.75,-4.5) rectangle (9.25,3.5); \fill[Teal]  (8.2,-2) rectangle (9.8,1);
  \end{scope}

  % ── 家族 5：多维 / 矩阵 ────────────────────────────────────────────────
  \node[family label,anchor=west] at (10.1,18.02) {05 · 多维与矩阵};

  % Parallel
  \node[scene card,anchor=north west] (sc-parallel) at (10.1,17.69) {};
  \node[scene caption] at (10.15,16.04)
    {{\bfseries Parallel}\\[0.5pt]\textcolor{Slate}{远征剖面}\\\textcolor{Slate}{五维野外条件}};
  % 小装饰：平行坐标（折线为填充细带、分层互不交叉，轴线只留上下短刻度）
  \begin{scope}[shift={(sc-parallel.center)},x=1mm,y=1mm]
    \foreach \x in {-9,-4.5,0,4.5,9} {
      \draw[white!35,line width=.5pt] (\x,3.2) -- (\x,3.9);
      \draw[white!35,line width=.5pt] (\x,-3.2) -- (\x,-3.9);
    }
    \fill[Teal,opacity=.85]
      (-9,2.525) -- (-4.5,1.725) -- (0,2.325) -- (4.5,1.925) -- (9,2.725)
      -- (9,2.475) -- (4.5,1.675) -- (0,2.075) -- (-4.5,1.475) -- (-9,2.275) -- cycle;
    \fill[Gold,opacity=.85]
      (-9,0.525) -- (-4.5,-0.275) -- (0,0.325) -- (4.5,-0.075) -- (9,0.725)
      -- (9,0.475) -- (4.5,-0.325) -- (0,0.075) -- (-4.5,-0.525) -- (-9,0.275) -- cycle;
    \fill[Coral,opacity=.85]
      (-9,-1.475) -- (-4.5,-2.275) -- (0,-1.675) -- (4.5,-2.475) -- (9,-1.275)
      -- (9,-1.525) -- (4.5,-2.725) -- (0,-1.925) -- (-4.5,-2.525) -- (-9,-1.725) -- cycle;
  \end{scope}

  % Heatmap
  \node[scene card,anchor=north west] (sc-heatmap) at (13.6,17.69) {};
  \node[scene caption] at (13.65,16.04)
    {{\bfseries Heatmap}\\[0.5pt]\textcolor{Slate}{服务依赖压力}\\\textcolor{Slate}{调用频率 × 延迟}};
  % 小装饰：热力矩阵（纯填充）
  \begin{scope}[shift={(sc-heatmap.center)},x=1mm,y=1mm]
    \foreach \i in {0,...,35} {
      \pgfmathtruncatemacro{\cx}{mod(\i,6)}
      \pgfmathtruncatemacro{\cy}{\i/6}
      \pgfmathtruncatemacro{\dist}{abs(\cx-\cy)}
      \ifnum\dist=0 \definecolor{tmpc2}{HTML}{1a2f4a} \else
      \ifnum\dist=1 \definecolor{tmpc2}{HTML}{2E6B8A} \else
      \ifnum\dist=2 \definecolor{tmpc2}{HTML}{2FA496} \else
      \ifnum\dist=3 \definecolor{tmpc2}{HTML}{54d6c6} \else
      \ifnum\dist=4 \definecolor{tmpc2}{HTML}{E8B94F} \else
      \definecolor{tmpc2}{HTML}{E2667B} \fi\fi\fi\fi\fi
      \fill[tmpc2] (\cx*2.3 - 5.75, \cy*2.3 - 5.75)
            rectangle (\cx*2.3 - 3.75, \cy*2.3 - 3.75);
    }
  \end{scope}

  % ── 家族 6：监控 / 剖面 ────────────────────────────────────────────────
  \node[family label,anchor=west] at (0,14.76) {06 · 监控与对比};

  % Gauge Cluster
  \node[scene card,anchor=north west] (sc-gauge) at (0,14.43) {};
  \node[scene caption] at (0.05,12.78)
    {{\bfseries Gauge Cluster}\\[0.5pt]\textcolor{Slate}{系统健康仪表}\\\textcolor{Slate}{容量/积压/命中率}};
  % 小装饰：四个仪表（底环与数值弧均为填充环形扇区，同心不同半径，互不重叠；
  % 用填充而非描边弧，避免被审计计为无标签连线）
  \begin{scope}[shift={(sc-gauge.center)},x=1mm,y=1mm]
    \fill[white!20] ([shift={(-10,0)}]180:2.76) arc (180:0:2.76)
      -- ([shift={(-10,0)}]0:2.44) arc (0:180:2.44) -- cycle;
    \fill[Teal] ([shift={(-10,0)}]180:1.96) arc (180:306:1.96)
      -- ([shift={(-10,0)}]306:1.64) arc (306:180:1.64) -- cycle;

    \fill[white!20] ([shift={(-4,0)}]180:2.76) arc (180:0:2.76)
      -- ([shift={(-4,0)}]0:2.44) arc (0:180:2.44) -- cycle;
    \fill[Gold] ([shift={(-4,0)}]180:1.96) arc (180:279:1.96)
      -- ([shift={(-4,0)}]279:1.64) arc (279:180:1.64) -- cycle;

    \fill[white!20] ([shift={(2,0)}]180:2.76) arc (180:0:2.76)
      -- ([shift={(2,0)}]0:2.44) arc (0:180:2.44) -- cycle;
    \fill[Coral] ([shift={(2,0)}]180:1.96) arc (180:333:1.96)
      -- ([shift={(2,0)}]333:1.64) arc (333:180:1.64) -- cycle;

    \fill[white!20] ([shift={(8,0)}]180:2.76) arc (180:0:2.76)
      -- ([shift={(8,0)}]0:2.44) arc (0:180:2.44) -- cycle;
    \fill[Violet] ([shift={(8,0)}]180:1.96) arc (180:351:1.96)
      -- ([shift={(8,0)}]351:1.64) arc (351:180:1.64) -- cycle;
  \end{scope}

  % Radar
  \node[scene card,anchor=north west] (sc-radar) at (3.3,14.43) {};
  \node[scene caption] at (3.35,12.78)
    {{\bfseries Radar}\\[0.5pt]\textcolor{Slate}{航天器韧性剖面}\\\textcolor{Slate}{三方案 × 六子系统}};
  % 小装饰：雷达剖面（两个半透明填充多边形，无描边）
  \begin{scope}[shift={(sc-radar.center)},x=1mm,y=1mm]
    \fill[Teal,opacity=.4]
      (90:5.5) -- (150:4.5) -- (210:5) -- (270:4) -- (330:5.5) -- (30:4.5) -- cycle;
    \fill[Gold,opacity=.4]
      (90:3.5) -- (150:5.5) -- (210:3) -- (270:5) -- (330:4) -- (30:5.5) -- cycle;
  \end{scope}

  % 项目核心文件（第三行右侧）
  \node[module,anchor=north west,minimum width=57mm,fill=Panel,
        text width=53mm,align=left,inner sep=2mm]
    (maints) at (6.8,14.43)
    {{\footnotesize\bfseries src/main.ts}\\[1pt]
     {\scriptsize 单文件场景注册}\\
     {\scriptsize 12 个 EChartsOption → 场景路由}\\
     {\scriptsize ?scene= 切换}};

  \node[module,anchor=north west,minimum width=57mm,fill=Panel,
        text width=53mm,align=left,inner sep=2mm]
    (capture) at (12.9,14.43)
    % \kern：pdftocairo 导出时 .mjs 后缀相邻字形 bbox 亚像素相切，svg 审计会误报
    {{\footnotesize\bfseries scripts/capture.m{\kern0.8pt}j{\kern0.8pt}s}\\[1pt]
     {\scriptsize Playwright 离线渲染}\\
     {\scriptsize 逐场景截图 + 交互测试}\\
     {\scriptsize RGBA 像素级校验}};

  % 分隔线
  \draw[boundary] (0,11.0) -- (19.1,11.0);

  % ════════════════════════════════════════════════════════════════════════
  %  第三段：构建与渲染管线（主路径）
  % ════════════════════════════════════════════════════════════════════════

  \node[kicker,anchor=west] at (0,10.62) {渲染管线 · Render Pipeline};
  \node[section title,anchor=west] at (0,10.02) {从 TypeScript 源到透明 PNG 全离线闭环};
  \node[section sub,anchor=north west,text width=68mm,align=left] at (0,9.6)
    {npm test = 类型检查 + 构建 + Playwright 渲染 + RGBA 质量断言，整个过程无外部网络请求。};

  % 信任边界线（本地回环）
  \draw[boundary] (13.6,10.62) -- (18.6,10.62);
  \node[edge label,anchor=east] at (13.35,10.62) {本地回环 · 无外网};

  % Vite server 模块
  \node[module,anchor=center,minimum width=34mm,fill=Panel,
        text width=31mm,align=center,inner sep=1.6mm]
    (vite) at (15.0,9.9)
    {{\footnotesize\bfseries Vite Dev Server}\\[0.5pt]
     {\scriptsize 127.0.0.1:4175}};

  % ── 六步主路径（编号与名称合一的路标 pill）──────────────────────────────
  \node[stage pill] (p1) at (2.0,8.2) {01 TS 源码};
  \node[stage pill] (p2) at (4.75,8.2) {02 类型检查};
  \node[stage pill] (p3) at (7.5,8.2) {03 Vite 构建};
  \node[stage pill] (p4) at (10.25,8.2) {04 本地服务};
  \node[stage pill] (p5) at (13.0,8.2) {05 浏览器渲染};
  \node[stage pill] (p6) at (15.75,8.2) {06 PNG 输出};

  \node[terminal,anchor=east] at (0.45,8.2) {main.ts};
  \node[terminal,anchor=west] at (17.2,8.2) {out/*.png};

  % 主色带：分段绘制，每段两端都附着于路标
  \begin{pgfonlayer}{edge layer}
    \draw[spine] (0.8,8.2) -- (p1.west);
    \draw[spine] (p1.east) -- (p2.west);
    \draw[spine] (p2.east) -- (p3.west);
    \draw[spine] (p3.east) -- (p4.west);
    \draw[spine] (p4.east) -- (p5.west);
    \draw[spine] (p5.east) -- (p6.west);
    \draw[spine] (p6.east) -- (16.9,8.2);
  \end{pgfonlayer}

  % 色带段标签（置于色带下方留白处，对齐各段间隙中心）
  \node[edge label] at (3.38,7.52) {注册 12 场景};
  \node[edge label] at (6.1,7.52) {tsc 通过};
  \node[edge label] at (8.9,7.52) {dist/};
  \node[edge label] at (14.4,7.52) {逐场景 PNG};
  \node[edge label] at (16.85,7.52) {产物};

  % Playwright / Chrome 模块（说明文字并入同一节点）
  \node[module,anchor=center,minimum width=36mm,fill=Panel,
        text width=33mm,align=center,inner sep=1.6mm]
    (pw) at (11.4,6.75)
    % \kern：pdftocairo 导出时相邻字形 bbox 亚像素相切（'+'/'C'、'f'/'t'），svg 审计会误报
    {{\footnotesize\bfseries Playwright\ {\kern0.8pt}+{\kern0.8pt}\ Chrome}\\[0.5pt]
     {\scriptsize headless · Swif{\kern0.9pt}tShader\\阻断外部网络}};

  % 调用 / 返回
  \begin{pgfonlayer}{edge layer}
    % Vite 往返（上方；call 挂在 p4 以绕开左侧大标题的横排文字区）
    \draw[call] (p4.north) -- (p4.north |- vite.west) -- (vite.west);
    \draw[reply] (vite.south) -- ++(0,-0.45) -| (p5.north);
    \node[edge label] at (11.7,9.6) {构建产物 dist/};
    \node[edge label,anchor=south] at (14.0,9.06) {HTTP 提供};

    % Playwright 往返（下方）
    \draw[call] (p4.south) -- (pw.north -| p4.south);
    \draw[reply] (pw.north -| p5.south) -- (p5.south);
    \node[edge label,anchor=west] at (10.35,7.6) {加载页面};
    \node[edge label,anchor=east] at (12.9,7.6) {截图结果};
  \end{pgfonlayer}

  % 分隔线
  \draw[boundary] (0,5.85) -- (19.1,5.85);

  % ════════════════════════════════════════════════════════════════════════
  %  第四段：三重质量门
  % ════════════════════════════════════════════════════════════════════════

  \node[kicker,anchor=west] at (0,5.5) {质量验证 · Quality Gates};
  \node[section title,anchor=west] at (0,4.95) {三重断言确保每个场景都「真的渲染出来了」};
  \node[section sub,anchor=west] at (0,4.48)
    {不依赖肉眼审查，每条断言都在 CI 中自动执行。};

  % 三个质量门卡片
  \begin{scope}[shift={(0,1.1)}]
    % Gate 1
    \node[module,anchor=north west,minimum width=58mm,minimum height=30mm,fill=Panel]
      (g1) at (0,3.0) {};
    \node[anchor=north west,font=\rmfamily\footnotesize\bfseries,text=Ink] at (0.3,2.72) {门 1 · 类型检查};
    \node[note,anchor=north west,align=left,text width=52mm] at (0.3,2.14)
      {tsc --noEmit\\校验 12 个 EChartsOption 类型\\scene 路由表完整性};
    \node[pill,anchor=south east,inner xsep=2.4mm,inner ysep=1.2mm] at (5.42,0.34) {静态};

    % Gate 2
    \node[module,anchor=north west,minimum width=58mm,minimum height=30mm,fill=Panel]
      (g2) at (6.6,3.0) {};
    \node[anchor=north west,font=\rmfamily\footnotesize\bfseries,text=Ink] at (6.9,2.72) {门 2 · 交互契约};
    \node[note,anchor=north west,align=left,text width=52mm] at (6.9,2.14)
      {等待 \_\_VIS\_READY\_\_ 标志\\注入 pointer 事件\\验证 \_\_INTERACTION\_COUNT\_\_ 递增};
    \node[pill,anchor=south east,inner xsep=2.4mm,inner ysep=1.2mm,fill=Gold!30,text=Gold!60!black] at (12.02,0.34) {运行时};

    % Gate 3
    \node[module,anchor=north west,minimum width=58mm,minimum height=30mm,fill=Panel]
      (g3) at (13.2,3.0) {};
    \node[anchor=north west,font=\rmfamily\footnotesize\bfseries,text=Ink] at (13.5,2.72) {门 3 · RGBA 像素断言};
    \node[note,anchor=north west,align=left,text width=52mm] at (13.5,2.14)
      {透明像素 > 8\%\\可见像素 > 3.5\%\\彩色像素 > 2500\\防止空白 / 纯黑输出};
    \node[pill,anchor=south east,inner xsep=2.4mm,inner ysep=1.2mm,fill=Coral!25,text=Coral!80!black] at (18.62,0.34) {像素级};
  \end{scope}

  % 验证流程箭头（门间依序通过）
  \begin{pgfonlayer}{edge layer}
    \draw[verify] (g1.east) -- (g2.west);
    \draw[verify] (g2.east) -- (g3.west);
  \end{pgfonlayer}
  \node[edge label,anchor=south] at (6.2,2.72) {通过};
  \node[edge label,anchor=south] at (12.8,2.72) {通过};

  % 分隔线
  \draw[boundary] (0,0.75) -- (19.1,0.75);

  % ════════════════════════════════════════════════════════════════════════
  %  第五段：技术栈与设计原则（底部安静区）
  % ════════════════════════════════════════════════════════════════════════

  \node[kicker,anchor=west] at (0,0.4) {技术栈 · Stack};

  % 左列：依赖
  \node[family label,anchor=west] at (0,-0.15) {运行时};
  \node[note,anchor=north west,align=left] at (0,-0.5)
    {echarts 6.1.0\\Canvas 渲染器\\透明背景设计};

  \node[family label,anchor=west] at (4.8,-0.15) {构建};
  \node[note,anchor=north west,align=left] at (4.8,-0.5)
    {TypeScript 5.9\\Vite 8.2\\ESM 模块};

  \node[family label,anchor=west] at (9.6,-0.15) {测试 / 渲染};
  \node[note,anchor=north west,align=left] at (9.6,-0.5)
    {Playwright + Chrome\\pngjs 像素分析\\离线网络隔离};

  \node[family label,anchor=west] at (14.4,-0.15) {设计原则};
  \node[note,anchor=north west,align=left] at (14.4,-0.5)
    {场景即文档\\目录驱动组织\\断言胜过肉眼};

  % 源注
  \node[note,anchor=west] at (0,-1.85)
    {数据来源：项目 README.md · catalog.json · src/main.ts · scripts/capture.mjs · package.json};

\end{tikzpicture}
\end{document}
